\section{USAMO 2022}
        \begin{enumerate}
        	\item (\textit{USAMO 2022}) Let $a$ and $b$ be positive integers. The cells of an $(a+b+1)\times (a+b+1)$ grid are colored amber and bronze such that there are at least $a^2+ab-b$ amber cells and at least $b^2+ab-a$ bronze cells. Prove that it is possible to choose $a$ amber cells and $b$ bronze cells such that no two of the $a+b$ chosen cells lie in the same row or column.


 


	\item (\textit{USAMO 2022}) Let $b\geq2$ and $w\geq2$ be fixed integers, and $n=b+w$. Given are $2b$ identical black rods and $2w$ identical white rods, each of side length 1.


 We assemble a regular $2n$-gon using these rods so that parallel sides are the same color. Then, a convex $2b$-gon $B$ is formed by translating the black rods, and a convex $2w$-gon $W$ is formed by translating the white rods. An example of one way of doing the assembly when $b=3$ and $w=2$ is shown below, as well as the resulting polygons $B$ and $W$.


 
                    \begin{figure}[H]
                        \centering
                        \includegraphics[width=0.5\textwidth]{Problem/USAMO_I_2022_2_Problem_Figure_1.png}
                    \end{figure}
                    


 Prove that the difference of the areas of $B$ and $W$ depends only on the numbers $b$ and $w$, and not on how the $2n$-gon was assembled.


 



\textbf{Solution 1}

First notice that the black rods and the white rods form polygons iff in the original $2n$-gon, if a side is a color $x$, then the side that is parallel to that side in the original $2n$-gon is also the color $x$. 

 We can prove that the difference in areas is only affected by the values of $b$ and $w$ by showing that for any valid arrangement of $2b$ rods and $2w$ rods, we may switch any two adjacent black and white rods(and their "parallel pairs"), and end up with the same area difference. 

 
                    \begin{figure}[H]
                        \centering
                        \includegraphics[width=0.5\textwidth]{Solution/USAMO_I_2022_2_Solution_Figure_1.png}
                    \end{figure}
                    

 In the figure above (click to expand), after the switch, we can see that after removing the mutually congruent parts, we are left with two parallelograms from each color. Let $x$, $\alpha{}$, $y$, and $\beta{}$ be defined as shown. Notice that if we angle chase, the sides of the other parallelogram are the same, but if the angles of the original $2n$-gon all have measure $2k$, the angles of the new parallelograms are $\alpha{}+180-2k$ and $\beta{}+180-2k$, as shown. We must prove that the differences between the areas are the same. 

 Using area formulas, the change in the difference of areas is $x\sin{\alpha{}}-y\sin{\beta{}}+y\sin{(2k-\beta{})}-x\sin{(2k-\alpha{})}$, which is equal to $x(\sin{k}(2\cos^2{k})-2\sin{k}\cos{k}\cos{\alpha{}})-y(\sin{k}(2\cos^2{k}))+y(\sin{k}\cos{k}\cos{\beta{}})$, or $2\cos{k}(x\sin{(\alpha{}-k)}+y\sin{(k-\beta{}}))$. Since $\cos{k}$ is not $0$ because $n\geq{}3$, we are left with proving that $x\sin{(\alpha{}-k)}+y\sin{(k-\beta{}})=0$.

 Now we rotate the polygon so that the vertex between the two sides that we switched is at the point $(0,0)$, the angle bisector of that vertex is $y=0$, and the black side is in the positive $y$-direction. Now think of all the sides as vectors, all pointing in the clockwise direction of the $2n$-gon.

 Notice the part labeled $a$ in the black polygons. We have that the vector labeled $x$ is really just the sum of all of the vectors in the part labeled $a$ - or all the vectors in the $2n$-gon that are in the positive $y$-direction excluding the one that was interchanged. Also notice that the angle of this vector $x$ has a signed angle of $\alpha{}-k$ with $y=0$ and has length $x$ - meaning that the vertical displacement of the vector $x$ from $y=0$ is equal to $x\sin{(\alpha{}-k)}$! Similarly, we get that the vertical displacement of the vector $y$ is equivalent to $y\sin{(k-\beta{}})$. 

 Adding these two together, we get that $x\sin{(\alpha{}-k)}+y\sin{(k-\beta{}})$ is simply the vertical displacement of the sum of the vectors $x$ and $y$. Since the sum of the vectors $x$ and $y$ is equivalent to the sum of the vectors in the positive half of the polygon minus the sum of the black vector that would be switched with the white vector(the leftmost vector in the positive half of the polygon) and the rightmost vector in the positive half(which is the parallel pair of the white vector that would be interchanged later), and we know that this sum happens to have a vertical displacement of $0$, along with the fact that the positive half of the polygon summed together also has a vertical displacement of $0$, we get that the total vertical displacement is $0$, meaning that $x\sin{(\alpha{}-k)}+y\sin{(k-\beta{}})=0$, and we are done.

 


 



\textbf{Solution 2}

Pick a pair of parallel sides of the regular $2n$-gon and flip their color. Let those two lines be the horizontal. The area of the polygon formed by rods of the original color decreases by a parallelogram with height equal to the sum of the vertical heights of the rods of that color divided by $2$ and base length $1$. Similarly, the area of the polygon formed by rods of the new color increases by the sum of the vertical heights of the rods of the new color divided by $2$. The sum of the heights of all the rod is equal to twice the height of the $2n$-gon, so the difference between the areas of $B$ and $W$ changes by the height of the $2n$-gon, which is fixed when $n$ is fixed.

 If we pair two pairs different-colored pairs of parallel sides of the $2n$-gon and flip the colors of both pairs, then $b$ and $w$ and the difference in the areas of $B$ and $W$ will remain unchanged. Thus, the difference in the areas of $B$ and $W$ depends only on $b$ and $w$.

 


	\item (\textit{USAMO 2022}) Let $\mathbb{R}_{>0}$ be the set of all positive real numbers. Find all functions $f:\mathbb{R}_{>0} \to \mathbb{R}_{>0}$ such that for all $x,y\in \mathbb{R}_{>0}$ we have

 \[f(x) = f(f(f(x)) + y) + f(xf(y)) f(x+y).\]


 



\textbf{Solution}

Denote the $n$th iteration of the function $f(x)$ as $f^{(n)}(x)$.

 First we prove that, for any $x$, $y>0$ such that $f(x)\leq f(y)$, we have $f^{(2)}(x)\geq y$. (1)

 Otherwise, $y>f^{(2)}(x)$, so $y-f^{(2)}(x)>0$; we call this value $c$, and since $c>0$, we can substitute $y$ with $c$ in the equation given in the problem, which gives 
 \[f(x)=f(f(f(x))+c)+f(xf(c))f(x+c)\] Recall that $c=y-f^{(2)}(x)$, which means $f(f(x))+c=y$; so $f(x)=f(y)+f(xf(c))f(x+c)$. However, since $f(x)$ is defined on $\mathbb{R}_{>0} \rightarrow \mathbb{R}_{>0}$, we know $f(xf(c))>0$, $f(x+c)>0$; so that $f(xf(c))f(x+c)>0$, and $f(x)=f(y)+f(xf(c))f(x+c)>f(y)$, which is a contradiction.

 
Now, since $f(x)\leq f(x)$, we know $f^{(2)}(x)\geq x$. Therefore, $f^{(4)}(x)=f^{(2)}(f^{(2)}(x))\geq f^{(2)}(x)$.

 Also, we know $f^{(3)}(x)=f^{(2)}(f(x))\geq f(x)$, and hence $f^{(5)}(x)=f^{(2)}(f^{(3)}(x)) \geq f^{(3)}(x)$.

 This means $f(x)\leq f^{(5)}(x)$, which is $f(x)\leq f(f^{(4)}(x))$; from (1) we know $f^{(2)}(x)\geq f^{(4)}(x)$. Since $f^{(2)}(x)\geq f^{(4)}(x)$, while we also obtained $f^{(4)}(x)\geq f^{(2)}(x)$, we know $f^{(2)}(x)=f^{(4)}(x)$.

 Then we try to prove that $f(x)$ is injective.

 If $f(x)$ is not injective, this means there exists some $a<b$ such that $f(a)=f(b)$.

 First consider substituting $x$ with $f^{(2)}(x)$ and $y$ with $a$, hence we obtain 
 \[f(f^{(2)}(x))=f(f(f(f^{(2)}(x)))+a)+f(f^{(2)}(x)f(a))f(f^{(2)}(x)+a)\] which means $f^{(3)}(x)=f(f^{(4)}(x)+a)+f(f^{(2)}(x)f(a))f(f^{(2)}(x)+a)=f(f^{(2)}(x)+a)+f(f^{(2)}(x)f(a))f(f^{(2)}(x)+a)$, and that obtains $f^{(3)}(x)=f(f^{(2)}(x)+a)(1+f(f^{(2)}(x)f(a)))$.

 Symmetrically, if we substitute $x$ with $f^{(2)}(x)$ and $y$ with $b$, we shall obtain $f^{(3)}(x)=f(f^{(2)}(x)+b)(1+f(f^{(2)}(x)f(b)))$, so that 
 \[f(f^{(2)}(x)+a)(1+f(f^{(2)}(x)f(a)))=f(f^{(2)}(x)+b)(1+f(f^{(2)}(x)f(b)))\]

 However, as $f(a)=f(b)$, we have $(1+f(f^{(2)}(x)f(a)))=(1+f(f^{(2)}(x)f(b)))$; then we know 
 \[f(f^{(2)}(x)+a)=f(f^{(2)}(x)+b).\]

 Now we substitute only $y$ with $a$, to obtain 
 \[f(x)=f(f(f(x))+a)+f(x(f(a))f(x+a)\] which is $f(x)=f(f^{(2)}(x)+a)+f(x(f(a))f(x+a)$. Symmetrically substituting only $y$ with $b$, we get $f(x)=f(f^{(2)}(x)+b)+f(x(f(b))f(x+b)$, so finally

 
 \[f(f^{(2)}(x)+a)+f(x(f(a))f(x+a)=f(f^{(2)}(x)+b)+f(x(f(b))f(x+b)\]

 But we know $f(f^{(2)}(x)+a)=f(f^{(2)}(x)+b)$, and as $f(a)=f(b)$, $f(x(f(a))=f(xf(b))$, so finally we know $f(x+a)=f(x+b)$.

 This means $f(x)$ is a periodic function with $b-a$ as a period. We can always take an arbitrarily large integer $n$, so that $f^{(2)}(a)<a+n(b-a)$.

 Recall the statement of (1), whose contrapositive is "if $f^{(2)}(x)<y$, we have $f(x)>f(y)$". By comparing $f^{(2)}(a)<a+n(b-a)$, we know $f(a)>f(a+n(b-a))=f(a)$, which is a contradiction!

 Hence, we finished the proof that $f(x)$ is injective.

 Then,as $f^{(4)}(x)=f^{(2)}(f^{(2)}(x))=f^{(2)}(x)$, we can conclude that $f^{(2)}(x)=x$.

 Recall the original equation, which $f(x)=f(f(f(x))+y)+f(xf(y))f(x+y)=f(f^{(2)}(x)+y)+f(xf(y))f(x+y)$, and it further means 
 \[f(x)=f(x+y)+f(xf(y))f(x+y)\].

 Substituting $x$ with $1$ gives $f(1)=f(1+y)+f(1\cdot f(y))f(1+y)=f(1+y)+f(f(y))f(1+y)=f(1+y)+yf(1+y)$, so $f(1)=f(1+y)(1+y)$ and thus $f(1+y)=\frac{f(1)}{1+y}$. Denote $f(1)$ as d, so we have proved that for all $z=1+y>1$, it is true that $f(z)=\frac{d}{z}$.

 For $z\leq 1$, we substitute $y$ with $f(z)$, which gives $f(x)=f(x+f(z))+f(xf(f(z)))f(x+f(z))=f(x+f(z))+f(xz)f(x+f(z))$. As the choice for $x$ is arbitrary, we can always take some $x>\frac{1}{z}$. Since $z\leq 1$, we know $x>1$, and $xz>\frac{1}{z}\cdot z=1$. So for the equation $f(x)=f(x+f(z))+f(xz)f(x+f(z))$, we have $x>1$, $x+f(z)>x>1$, and $xz>1$. Then we can write 

 
 \[\frac{d}{x}=\frac{d}{x+f(z)}+\frac{d}{xz}\cdot\frac{d}{x+f(z)}\]

 Multiplying $\frac{xz(x+f(z))}{d}$ on both sides give $z(x+f(z))=xz+d$, so $f(z)=\frac{d}{z}$.

 We have finally proved that, for all $z$, it is true that $f(z)=\frac{d}{z}$; and we thus conclude our final answer, that 
 \[\boxed{f(x)=\frac{d}{x}}\] where $d\in\mathbb{R}_{>0}$. 

 


	\item (\textit{USAMO 2022}) Find all pairs of primes $(p, q)$ for which $p-q$ and $pq-q$ are both perfect squares.


 



\textbf{Solution 1}

Since $q(p-1)$ is a perfect square and $q$ is prime, we should have $p - 1 = qb^2$ for some positive integer $b$. Let $a^2 = p - q$. Therefore, $q = p - a^2$, and substituting that into the $p - 1 = qb^2$ and solving for $p$ gives
 \[p = \frac{a^2b^2 - 1}{b^2 - 1} = \frac{(ab - 1)(ab + 1)}{b^2 - 1}.\]Notice that we also have
 \[p = \frac{a^2b^2 - 1}{b^2 - 1} = a^2 + \frac{a^2 - 1}{b^2 - 1}\]and so $b^2 - 1 | a^2 - 1$. We run through the cases

 
                \begin{itemize}
                    \item $a = 1$: Then $p - q = 1$ so $(p, q) = (3, 2)$, which works.\item $a = b$: This means $p = a^2 + 1$, so $q = 1$, a contradiction.\item $a > b$: This means that $b^2 - 1 < ab - 1$. Since $b^2 - 1$ can be split up into two factors $F_1, F_2$ such that $F_1 | ab - 1$ and $F_2 | ab + 1$, we get
                \end{itemize}


                
 \[p = \frac{ab - 1}{F_1} \cdot \frac{ab + 1}{F_2}\]and each factor is greater than $1$, contradicting the primality of $p$.

 Thus, the only solution is $(p, q) = (3, 2)$.

 



\textbf{Solution 2}

Let $p-q = a^2$, $pq - q = b^2$, where $a, b$ are positive integers. $b^2 - a^2 = pq - q - (p-q) = pq -p$. So,
 \[b^2 - a^2 = p(q-1) \tag{1}\]

 $\bullet$ For $q=2$, $p = b^2 - a^2 = (b-a)(b+a)$. Then $b-a=1$ and $b+a=p$. $a=\dfrac{p-1}{2}$ and $p-q = a^2$. Thus, $p - 2 = \left( \dfrac{p-1}{2} \right)^2 \implies p^2 - 6p + 9 = 0$ and we find $p=3$. Hence $(p,q) = (3,2)$.

 
$\bullet$ For $q=4k+3$, ($k\geq 0$ integer), by $(1)$, $p(4k+2) = b^2 - a^2$. Let's examine in $\mod 4$, $b^2 - a^2 \equiv 2 \pmod{4}$. But we know that $b^2 - a^2 \equiv 0, 1 \text{ or } 3 \pmod{4}$. This is a contradiction and no solution for $q = 4k + 3$.

 
$\bullet$ For $q=4k+1$, ($k > 0$ integer), by $(1)$, $p(4k) = b^2 - a^2$. Let $k=m\cdot n$, where $m\geq n \geq 1$ and $m, n$ are integers. Since $p>q$, we see $p>4k$. Thus, by $(1)$, $(b-a)(b+a) = 4p\cdot m \cdot n$. $b-a$ and $b+a$ are same parity and  $4p\cdot m \cdot n$ is even integer. So, $b-a$ and $b+a$ are both even integers. Therefore,

 $\left\{ \begin{array}{rcr} b+a =  & 2pn \\ b-a =  & 2m  \end{array} \right.$or $\left\{ \begin{array}{rcr} b+a =  & 2pm \\ b-a =  & 2n  \end{array} \right.$Therefore, $a=pn - m$ or $a = pm - n$. For each case, $p-q = p - 4mn - 1 < a$. But $p-q = a^2$, this gives a contradiction. No solution for $q = 4k + 1$. 

 
We conclude that the only solution is $(p,q) = (3,2)$.

 (Lokman GÖKÇE)

 



\textbf{Solution 3}

We claim that the only solution is $(p,q) = (3,2)$. First, this clearly works, since $3-2=1^2$ and $3\cdot2-2=2^2$.

 Start as in Solution 2 - assign positive integers $m, n,$ such that $p-q=m^2$ and $pq-q=n^2.$ Then, $pq-p=n^2-m^2,$ or $p(q-1)=(n-m)(n+m).$

 First, $p=q$ is impossible, since $p^2-p = p(p-1)$ is not a perfect square. Now notice that $m<p$ and $n<\max(p,q-1).$ But $p-q>0,$ so $p\geq q-1.$ Now, $m, n < p.$

 Therefore, $n-m$ can't ever be a factor of $p,$ and $n+m<2p,$ so $n+m=p.$ Then, $n-m=q-1,$ so $p-q=2m-1.$ But $p-q=m^2,$ so $m^2=2m-1,$ or $m=1.$ Therefore, one of $p$ and $q$ must be $2$ (and $1$ is not a prime).

 Thus, we conclude that our claim is true.

 (grape1984)

 


	\item (\textit{USAMO 2022}) A function $f: \mathbb{R}\to \mathbb{R}$ is \textit{essentially increasing} if $f(s)\leq f(t)$ holds whenever $s\leq t$ are real numbers such that $f(s)\neq 0$ and $f(t)\neq 0$.


 Find the smallest integer $k$ such that for any 2022 real numbers $x_1,x_2,\ldots , x_{2022},$ there exist $k$ essentially increasing functions $f_1,\ldots, f_k$ such that
 \[f_1(n) + f_2(n) + \cdots + f_k(n) = x_n\qquad \text{for every } n= 1,2,\ldots 2022.\]


 



\textbf{Solution}


 \[\textbf{Answer: } k = 1011.\]

 
 \[\textbf{Explanation:}\]

 
 \[\text{A function } f:\mathbb{R}\to\mathbb{R} \text{ is essentially increasing if } f(s)\le f(t) \text{ whenever } s\le t \text{ and } f(s)\ne 0,\, f(t)\ne 0.\]

 
 \[\text{We want the smallest } k \text{ such that for any } x_1,\dots,x_{2022}, \text{ we can write }\]

 
 \[f_1(n)+f_2(n)+\cdots+f_k(n)=x_n \quad \text{for all } n.\]

 
 \[\textbf{Lower bound:}\]

 
 \[\text{Consider } x_1 > x_2 > \cdots > x_{2022}.\]

 
 \[\text{Each essentially increasing function can contribute nonzero values}\]

 
 \[\text{on at most 2 positions (otherwise monotonicity is violated).}\]

 
 \[\Rightarrow k \ge \left\lceil \frac{2022}{2} \right\rceil = 1011.\]

 
 \[\textbf{Upper bound:}\]

 
 \[\text{Group indices as } (1,2), (3,4), \dots, (2021,2022).\]

 
 \[\text{For each pair, define one function that matches } x_{2i-1}, x_{2i}\]

 
 \[\text{and is zero elsewhere. This satisfies the condition.}\]

 
 \[\Rightarrow k \le 1011.\]

 


	\item (\textit{USAMO 2022}) There are 2022 users on a social network called Mathbook, and some of them are Mathbook-friends. (On Mathbook, friendship is always mutual and permanent.)


 Starting now, Mathbook will only allow a new friendship to be formed between two users if they have \textit{at least two} friends in common. What is the minimum number of friendships that must already exist so that every user could eventually become friends with every other user?


 



\textbf{Solution 1}

(but 4044 seems to be the wrong answer?)To answer this question, we need to consider how friendships on Mathbook are formed. If two users have at least two friends in common, then they will be able to become friends with each other. This means that the minimum number of friendships that must already exist in order for every user to eventually become friends with every other user is the smallest number of friendships that guarantees that each pair of users has at least two friends in common.

 One way to guarantee that each pair of users has at least two friends in common is to create a complete graph, where each user is friends with every other user. In this case, each pair of users has exactly $2022 - 2 = 2020$ friends in common. However, this is not the minimum number of friendships required, since some pairs of users may have more than two friends in common without forming a complete graph.

 To find the minimum number of friendships required, we can use the fact that a friendship between two users implies that they share all of their mutual friends. In other words, if two users are friends, then they must have at least two friends in common. This means that if we want to guarantee that each pair of users has at least two friends in common, we can do so by ensuring that each user has at least two friends.

 Therefore, the minimum number of friendships that must already exist in order for every user to eventually become friends with every other user is $2 \cdot 2022 = \boxed{4044}$.

 


\end{enumerate}